%!TEX root=../mythesis.tex
% Chapter 3 -- Literature Review

\chapter{Literature Review} % Main chapter title
\chaptermark{Literature Review}
\label{ch:literature_review}

This chapter surveys the prior work on which the three technical chapters build. With the financial and reinforcement-learning preliminaries of \cref{ch:preliminaries} in hand, \sref{sec:bg_rl4qt} reviews reinforcement learning for quantitative trading organised by task---single-asset and intraday trading, high-frequency trading, portfolio management, order execution, strategy routing, and formulaic alpha mining---before turning to the cross-cutting question of uncertainty and non-stationarity and closing with further related work. \sref{sec:bg_gaps} then synthesises the open challenges that motivate \cref{ch:earnhft}, \cref{ch:fineft}, and \cref{ch:alphaqd}.

%----------------------------------------------------------------------------------------
\section[Reinforcement Learning for Quantitative Trading]{Reinforcement Learning for\\Quantitative Trading}
\label{sec:bg_rl4qt}

Reinforcement learning has been applied across the entire quantitative-trading stack, and we survey that body of work organised by task. For each task we note the dominant design and the assumption it leaves unaddressed, which the open challenges of \sref{sec:bg_gaps} then collect. Open frameworks and benchmarks have standardised much of this research---FinRL~\cite{liu2021finrl}, TradeMaster~\cite{sun2024trademaster}, Qlib~\cite{yang2020qlib}, and the PRUDEX evaluation compass~\cite{sun2023prudex}---and the studies in this thesis follow their evaluation conventions where applicable.

\subsection{Single-Asset and Intraday Trading}
\label{sec:bg_intraday}

The earliest and largest body of RL-for-trading work targets a single asset at the daily-to-intraday scale, where the agent maps a window of price and technical features to a discrete trading action. Foundational studies established deep RL as a viable trader on this task~\cite{zhang2019deep,theate2021application,li2019deep}, and later work improved robustness and sample efficiency through imitation of profitable behaviour~\cite{liu2020adaptive,ding2018investor}, multimodal or convolutional state encoders~\cite{shin2019deep}, reward shaping for sparse-reward environments~\cite{takara4411793deep}, and multi-objective formulations for intraday trading of stock-index futures~\cite{si2017multi}. DeepScalper~\cite{sun2022deepscalper} adds a hindsight bonus and an auxiliary price-prediction task to improve generalisation in intraday trading and remains a strong representative of the single-agent design. CDQNRP~\cite{zhu2022quantitative}, sometimes reused as a baseline in HFT studies, evaluates a perturbed convolutional DQN at daily frequency and belongs to this adjacent single-asset group rather than to HFT proper. The common thread is a single, monolithic policy: effective at its target frequency, but not by itself designed to address the horizon growth of second-level trading or the regime non-stationarity that erodes one policy over long deployments.

\subsection{High-Frequency Trading}
\label{sec:bg_hft}

For high-frequency trading specifically, DRA~\cite{briola2021deep} couples an LSTM~\cite{hochreiter1997long} state encoder with PPO, and asynchronous deep double-duelling Q-learning has been applied to rapid trading-signal decisions~\cite{nagy2023asynchronous}. A common non-learned alternative is rule-based technical analysis, built on signals such as MACD~\cite{appel2005technical} and the five-level order-book volume-imbalance rule described in \sref{sec:bg_instruments}. EarnHFT addresses second-level position control through hierarchy; FineFT is a separate five-minute leveraged-futures study rather than an HFT method.

\subsection{Portfolio Management and Asset Allocation}
\label{sec:bg_portfolio}

Beyond single-asset timing, a substantial line applies RL to portfolio management, where the action is a continuous allocation across many assets. Early frameworks cast allocation as a deterministic policy-gradient problem over price tensors~\cite{jiang2017deep,jiang2017cryptocurrency}; subsequent work added explicit risk--return balancing~\cite{wang2021deeptrader}, ensemble strategies that switch among actor--critic learners~\cite{yang2020deep}, hierarchical decomposition of selection, allocation, and commission-aware trade execution~\cite{gao2021framework,zha2022hierarchical,wang2021commission}, and maskable representations that respect tradability constraints~\cite{zhang2024reinforcement}. Margin Trader~\cite{gu2023margin} extends the setting to leverage and short selling---the regime that \cref{ch:fineft} studies. Portfolio management is not the subject of this thesis, but it shares the central difficulty of non-stationarity and motivates much of the ensemble and hierarchical machinery reused here.

\subsection{Order Execution}
\label{sec:bg_execution}

A complementary task is optimal order execution---acquiring or liquidating a large position while minimising market impact. Recent RL work on execution stresses generalisation across instruments and regimes~\cite{zhang2023towards}, and the adjacent microstructure task of automated market making has been tackled with imitative RL over predictive state representations~\cite{niu2023imm}. Execution is the task in which EarnHFT's historical-replay assumption---that simulated actions do not alter the recorded price path (\sref{sec:earnhft_stage1})---is most visibly an approximation, and the literature here is correspondingly attentive to market impact.

\subsection{Hierarchical and Routing Approaches}
\label{sec:bg_routing}

A growing line of work does not commit to one policy but routes among several. MetaTrader~\cite{niu2022metatrader} learns a router that picks the most suitable strategy for the current market, and mixture-of-experts designs train diversified traders and combine them~\cite{sun2023mastering}. Winnow~\cite{littlestone1988winnow} is originally a mistake-driven linear-threshold learning algorithm; the trading experiments later use a performance-based adaptation under that name, so it is not treated here as a generic theorem for expert weighting. EarnHFT~\cite{qin2024earnhft} and MacroHFT~\cite{zong2024macrohft} construct hierarchical MDPs for high-frequency trading and train the high-level router \emph{by reinforcement learning on historical data}; in EarnHFT the hindsight Q supervision applies only to the second-level specialists, and the router itself is trained without a Q-teacher (\cref{ch:earnhft}). In both systems the router is fitted offline and then frozen for deployment rather than updated online.

\subsection{Formulaic Alpha Mining}
\label{sec:bg_alpha}

Formulaic alpha mining seeks short symbolic expressions whose cross-sectional values predict returns. Hand-crafted catalogues such as Alpha101~\cite{kakushadze2016101} gave way to automated search. Genetic programming~\cite{stephens2015gplearn} and its hierarchical and warm-started variants~\cite{zhang2020autoalpha,ren2024warmstart} evolve expression trees by mutation and crossover. The broader quality-diversity lineage seeks a repertoire rather than one optimum~\cite{pugh2016quality}; MAP-Elites discretises a behaviour space and retains high-quality occupants of its cells~\cite{mouret2015mapelites} (see also \sref{sec:bg_qd}). AutoAlpha is especially relevant because its PCA-QD component is an earlier use of these ideas in alpha mining; AlphaQD's novelty is therefore the separate return-source MAP-Elites archive and its separation from the generator reward, not QD in the domain per se. AlphaEvolve~\cite{cui2021alphaevolve} is another learning framework for discovering novel alphas and is reviewed here even though it is absent from the unified experimental table. Deep symbolic optimisation (DSO)~\cite{petersen2021dso} trains a recurrent generator against a standalone expression score, providing a counterexample to the claim that every RL-style symbolic generator must use a pool-coupled reward. The pool-based reinforcement-learning line includes AlphaGen~\cite{yu2023alphagen}, which casts mining as a PPO problem whose terminal reward is the marginal contribution of a new alpha to a running pool. AlphaQCM~\cite{zhu2025alphaqcm} adds a distributional value head and a variance bonus; RiskMiner~\cite{ren2024riskminer} pairs risk-seeking Monte-Carlo tree search with a pool reward; Alpha\textsuperscript{2}~\cite{xu2024alpha2} searches logical formulas; AlphaForge~\cite{shi2025alphaforge} replaces the RL search with a generator--predictor surrogate and a dynamic re-combination stage; and AlphaSAGE~\cite{chen2025alphasage} encodes expressions as typed graphs and trains a GFlowNet against a dense, multi-faceted reward. RiskMiner and Alpha\textsuperscript{2} are contextual methods rather than reported baselines here; their omission limits the coverage of the empirical comparison. A parallel and very recent line replaces the search policy with a large language model~\cite{wang2025alphagpt,tang2025alphaagent,li2024fama}. As \cref{ch:alphaqd} details, AlphaQD addresses the moving target in this influential pool-based family; it does not claim that every symbolic RL objective is pool-coupled.

\subsection{Uncertainty and Non-Stationary RL}
\label{sec:bg_uncertainty}

The leverage setting of \cref{ch:fineft} foregrounds two kinds of uncertainty. \emph{Aleatoric} uncertainty is the irreducible noise of the environment. Leverage does not change the underlying transition uncertainty, but it magnifies that noise in account returns and rewards; one risk-aware quantile construction evaluated in automated driving~\cite{bernhard2019addressing} and adapted to the trading setting in \cref{ch:fineft} targets the resulting lower-tail risk. Ensemble value methods address a different axis: Averaged-DQN reduces target variance~\cite{anschel2017averaged}, while SUNRISE uses ensemble disagreement in its sample weighting and action selection~\cite{lee2021sunrise}. \emph{Epistemic} uncertainty is the model's ignorance about states it has not seen and motivates OOD detection. A VAE~\cite{kingma2013auto} can supply a reconstruction-based distribution-familiarity score, although deep generative scores can be misleading under distribution shift~\cite{zhang2021understanding}; the score is a heuristic rather than calibrated epistemic uncertainty (see \sref{sec:bg_vae}). Modelling slowly drifting environments as a dynamic-parameter MDP~\cite{xie2020deep} provides the formal setting in which a short window can be treated as approximately stationary. FineFT combines selective ensemble updates with one such VAE score per inferred dynamic.

\subsection{Further Related Work}
\label{sec:bg_further}

Beyond the task-specific lines above, four strands of prior work help situate the contributions of this thesis. Moody and Saffell~\cite{moody2001learning} provide a foundational treatment of reinforcement learning for trading, training a recurrent trader by direct reinforcement on risk-adjusted performance rather than through an intermediate value function. Nevmyvaka, Feng, and Kearns~\cite{nevmyvaka2006reinforcement} give the classical large-scale study of reinforcement learning for optimised trade execution on limit-order-book data, the setting in which the exogenous-price approximation exploited by \cref{ch:earnhft} is most commonly invoked. The Q-teacher of \cref{ch:earnhft} is related to deep Q-learning from demonstrations~\cite{hester2018dqfd}, which shapes value learning with a supervised loss on expert data; the Q-teacher differs in that dynamic programming computes sample-path hindsight continuation values under the exogenous-replay assumption rather than using externally supplied demonstrations. These are trajectory-conditioned training signals, not optimal action values of the unknown stochastic market. Finally, Harvey, Liu, and Zhu~\cite{harvey2016cross} document how multiple testing inflates apparent performance in cross-sectional return research---a winner's-curse concern that applies directly to alpha mining, where many candidate formulas are evaluated on the same data (\cref{ch:alphaqd}).

%----------------------------------------------------------------------------------------
\section{Open Challenges}
\label{sec:bg_gaps}

This survey exposes three gaps that the technical chapters close.

\begin{itemize}
\item \textbf{Gap I (high-frequency efficiency and adaptivity).} Prior single-policy methods do not jointly address second-level horizon length and changing trends. A method is needed that improves learning efficiency at second scale and can route among differentiated policies. This motivates EarnHFT (\cref{ch:earnhft}).

\item \textbf{Gap II (risk and distribution awareness under leverage).} Methods designed for low-leverage markets do not necessarily control the magnified P\&L consequences of leverage or flag unfamiliar states, leaving them exposed during abrupt shifts. This motivates FineFT (\cref{ch:fineft}).

\item \textbf{Gap III (reward stationarity and selection in alpha mining).} A dominant family of pool-based RL alpha miners couples the reward to an evolving pool, training the policy against a moving target and complicating archive-conditional selection analysis. This motivates AlphaQD (\cref{ch:alphaqd}).
\end{itemize}

A thread common to all three gaps is the placement of the aggregation mechanism. EarnHFT learns it, FineFT derives it heuristically from generative uncertainty estimates, and AlphaQD removes it from the reward entirely---a progression that the remainder of the thesis develops.
